% Ортографическая проекция плоскости орбиты: (x,y,0) -> (x,y/2).
% Её нормаль проецируется в (0,sqrt(3)/2); a=25, e=0.3, nu=-55 deg.
\begin{tikzpicture}[course diagram,x=1mm,y=1mm]
  \path[use as bounding box] (-36,-21) rectangle (32,23);
  \pgfmathsetmacro{\orbitradius}{22.75/(1+.3*cos(-55))}
  \coordinate (F) at (0,0);
  \coordinate (O) at ({\orbitradius*cos(-55)},{.5*\orbitradius*sin(-55)});
  \coordinate (Raxis) at ($(O)+({18*cos(-55)},{9*sin(-55)})$);
  \coordinate (Taxis) at ($(O)+({-18*sin(-55)},{9*cos(-55)})$);
  \coordinate (Zaxis) at ($(O)+(0,{25*sqrt(3)/2})$);

  \fill[coursetint] (-7.5,0) ellipse[x radius=25,y radius=11.92424];
  \draw[courseaccent,line width=.8pt,dashed,domain=0:180,samples=65,variable=\t]
    plot ({-7.5+25*cos(\t)},{11.92424*sin(\t)});
  \draw[orbit line,domain=180:360,samples=65,variable=\t]
    plot ({-7.5+25*cos(\t)},{11.92424*sin(\t)});
  \draw[orbit line,->,domain=228:253,samples=20,variable=\t]
    plot ({-7.5+25*cos(\t)},{11.92424*sin(\t)});
  \node[text=coursemuted,align=center] at (-17,2)
    {Плоскость\\орбиты};
  \node[text=coursemuted,anchor=north] at (-18,-13) {Орбита КА};

  \draw[vector line] (F)--node[pos=.48,above right,inner sep=.7mm] {$\mathbf R$}(O);
  \fill[courseaccent] (F) circle[radius=1.5];
  \node[anchor=south,inner sep=1.5mm] at (F) {Земля};

  \draw[courseaccent,line width=1pt,->] (O)--(Raxis);
  \draw[courseaccent,line width=1pt,->] (O)--(Taxis);
  \draw[courseaccent,line width=1pt,->] (O)--(Zaxis);
  \node[text=courseaccent,anchor=north west,inner sep=.7mm] at (Raxis) {$\mathbf e_r$};
  \node[text=courseaccent,anchor=west,inner sep=.7mm] at (Taxis) {$\mathbf e_t$};
  \node[text=courseaccent,anchor=south,inner sep=1mm] at (Zaxis) {$\mathbf e_z$};
  \fill[courseink] (O) circle[radius=1];
  \draw[aux line] (O)--(7,-15.5);
  \node[anchor=north,inner sep=1mm] at (7,-15.5) {$O$ (КА)};
\end{tikzpicture}
